\begin{abstract}
传统辐射源定位先估计站间到达时间差，再由时差求解位置。多个未知波形占用同一频带时，强源会干扰弱源的时延提取与跨站关联。直接位置估计（DPD）利用多站混合观测直接拟合源位置，但高效求解仍面临两个问题：粗搜索可能过早舍弃弱源的有效候选，逐源局部更新又可能因位置与波形相互耦合而迭代缓慢。本文针对同步单天线站的未知宽带双源定位，采用已知自由空间相对幅度与基带时延模型，将未知波形以最小二乘消元，得到统一的位置目标。在此基础上提出多候选联合精修AP（MCJ-AP）：保留少量条件候选，分别同时调整两源位置，再按最终拟合残差选择结果。多候选用于减少错误区域选择，联合精修用于降低候选处理成本。另构造波形消元块NOMP，在局部位置更新中同步重估信号系数，并通过曲率分析解释其计算机制。在20个源站布局、六档功率差的2400条观测中，五候选MCJ-AP将双源20\,m成功率由单候选的95\%提高至99\%，平均求解时间为0.718\,s。相同五候选、相同停止精度的完整配对进一步表明，联合精修保持99\%成功率，将平均耗时从1.084\,s降至0.744\,s，减少31.36\%。共同240条观测上的八类方法比较表明，该流程兼顾了多源恢复能力与亚秒级平均计算成本。
\end{abstract}
\noindent\textbf{关键词：}直接位置估计；强弱重叠源；交替投影；块NOMP；可分离非线性最小二乘
